\section{Introduction}

When a trading agent that reads financial news improves its Sharpe ratio, the improvement can come from selecting different assets, holding a different fraction of capital in the market, or both. This ambiguity matters when comparing a news-driven agent with a fully invested reference. An explicit volatility rule and a budget-matched comparison let us examine the contribution of relative asset weights under the same invested fraction.

Large language models (LLMs) can reason over financial text, and prior work has studied whether sentiment extracted from news is associated with subsequent returns~\cite{lopezlira2023chatgpt}. FinMem and TradingGPT use memory to organise information for trading decisions~\cite{yu2024finmem,li2023tradinggpt}. FinCon combines multi-agent decision making with verbal reinforcement~\cite{yu2024fincon}, FinAgent augments its policy with tools~\cite{zhang2024finagent}, and TradingAgents assigns specialised roles to agents~\cite{xiao2024tradingagents}. FLAG-Trader fine-tunes an LLM trading policy through reinforcement learning~\cite{wang2025flagtrader}. When text interpretation feeds directly into a trading policy, changes in the model can affect both asset selection and portfolio exposure, making it difficult to attribute performance to information extracted from news. Benchmark studies document sensitivity to model choice and market conditions~\cite{li2024investorbench,li2025finsaber}. With date-only news, an article may describe a price move that has already occurred. A separate concern is look-ahead through information seen during pretraining~\cite{glasserman2023lookahead,lopezlira2025memorization}.

We propose \textbf{VolTilt}, a long-only allocator that assigns portfolio exposure to an explicit volatility rule and uses an LLM to inform relative asset weights (Section~\ref{sec:method}). The LLM extracts event records under a fixed schema. We retain articles that focus on the asset and report specific figures, smooth their impact scores, and use cross-sectional ranks to form a bounded tilt around equal weights. A common volatility budget scales the allocation using recent realised returns on the tilted portfolio~\cite{moreira2017volmanaged,harvey2018impact}. The text signal can therefore affect the volatility estimate through portfolio composition, while the scaling rule is fixed. Disabling the tilt gives a price-only comparison under the same rule. Holding the resulting budget fixed then lets us compare relative weights at the same daily invested fraction. We evaluate the allocator on three stock and cryptocurrency settings from the FINSABER and InvestorBench benchmarks against a fully invested equal-weight reference, rule-based strategies, and reimplementations of published LLM trading agents~\cite{li2025finsaber,li2024investorbench}.

\noindent\textbf{Our contributions are as follows.}
\begin{itemize}
  \item We develop a long-only allocation method that combines portfolio-level volatility scaling with a bounded rank tilt derived from structured financial news. The fixed event schema allows us to substitute LLM backbones while holding the scoring and allocation rules fixed.
  \item We assess the incremental value of the LLM signal using a price-only ablation, matched daily budgets, shifted-score and article-count controls, and tests of forward return associations under alternative news cutoffs (Section~\ref{sec:experiments}).
  \item We find that volatility scaling improves Sharpe and lowers drawdown relative to equal weight in the FINSABER stock window. The matched-budget comparisons across labellers offer limited evidence that the news tilt improves performance in the two stock settings.
\end{itemize}
